% =====================================================================
%  Modelo de consentimiento informado (PFC-ISR401-REPCOMP-A01, v1.0)
%  Se incluye como anexo desde main.tex. El ejemplar FIRMADO nunca se
%  sube al repositorio: va a la carpeta institucional privada y se
%  registra en evidencias/MANIFIESTO.md con su huella SHA-256.
% =====================================================================

\begin{center}
  {\small\bfseries FACULTAD DE CIENCIAS DE LA COMPUTACI\'ON $\cdot$ CARRERA DE INGENIER\'IA DE SOFTWARE\par}
  \vspace{0.1cm}
  {\large\bfseries FORMULARIO DE CONSENTIMIENTO INFORMADO\par}
  {\small Participaci\'on voluntaria en el levantamiento de requisitos de un proyecto acad\'emico\par}
\end{center}

\begin{tabularx}{\linewidth}{@{}L{4cm}Y@{}}
  \toprule
  \textbf{C\'odigo del documento} & PFC-ISR401-REPCOMP-A01 \quad Versi\'on 1.0 (16/11/2026) \\
  \textbf{Asignatura} & Ingenier\'ia de Requisitos (ISR-401) $\cdot$ 4.\textsuperscript{o} nivel $\cdot$ Paralelo A $\cdot$ Per\'iodo 2026--2027 \\
  \textbf{Proyecto} & Sistema de gesti\'on de reparaci\'on de equipos de c\'omputo \\
  \textbf{Organizaci\'on estudiada} & Taller de reparaci\'on de equipos de c\'omputo (nombre: JC\textbar{}cell), cant\'on Quevedo, provincia de Los R\'ios \\
  \textbf{Docente responsable} & Ing. Gleiston Cicer\'on Guerrero Ulloa, PhD $\cdot$ \texttt{gguerrero@uteq.edu.ec} \\
  \textbf{Contacto del equipo} & Basurto Anchaguano Jordan Steven $\cdot$ \texttt{jbasurtoa@uteq.edu.ec} \\
  \textbf{C\'odigo del participante} & \textbf{REPCOMP-P01} \\
  \bottomrule
\end{tabularx}

\medskip
Usted ha sido invitado o invitada a participar en una actividad acad\'emica de la Universidad
T\'ecnica Estatal de Quevedo (UTEQ). Antes de decidir, lea con calma este documento o pida que se
lo lean en voz alta. Puede hacer todas las preguntas que necesite. Participar es opcional: si no
desea hacerlo, no tiene que dar ninguna explicaci\'on.

\subsection*{1. ?`Para qu\'e es este estudio?}
Un grupo de estudiantes de Ingenier\'ia de Software est\'a estudiando c\'omo se trabaja en un taller
de reparaci\'on de equipos de c\'omputo, en particular en la recepci\'on y registro de equipos, el
diagn\'ostico y la cotizaci\'on, el seguimiento de las reparaciones, el control de repuestos e
inventario, la comunicaci\'on con los clientes y la facturaci\'on y entrega, con el fin de
identificar, analizar y documentar los requisitos de un sistema inform\'atico (Sistema de gesti\'on
de reparaci\'on de equipos de c\'omputo). La informaci\'on se usar\'a como evidencia para elaborar
documentos de ingenier\'ia de requisitos: listas de requisitos, casos de uso, modelos, prototipos y
registros de validaci\'on. Se trata de una propuesta acad\'emica: su participaci\'on no implica la
compra, instalaci\'on ni entrega de ning\'un sistema a la organizaci\'on.

\subsection*{2. ?`Qu\'e tendr\'ia que hacer?}
Seg\'un su rol (responsable del taller), se le pedir\'a una o varias de estas actividades, que el
equipo le indicar\'a antes de empezar:
\begin{itemize}[nosep]
  \item Una entrevista sobre sus tareas diarias, los problemas que encuentra y las mejoras que
        considera necesarias (aproximadamente 15 a 45 minutos).
  \item Responder a un formulario breve ($\square$ impreso o $\square$ en l\'inea; marque el que prefiera).
  \item Permitir que un integrante del equipo observe su trabajo sin intervenir (observaci\'on no participante).
  \item Revisar las interfaces de los prototipos del posible sistema y opinar si reflejan su trabajo
        (sesi\'on de validaci\'on).
\end{itemize}
Lo que usted diga se registrar\'a con notas escritas. El audio, el video y las fotograf\'ias solo se
aplicar\'an si usted lo autoriza en la Secci\'on 7. Negarse a ser grabado no le impide participar;
sin embargo, estos registros son de gran importancia en nuestro trabajo.

\textbf{No se le pedir\'a:} usuarios ni contrase\~nas, ni contenido personal de los equipos de los
clientes (archivos, fotos, correos, conversaciones); tampoco datos bancarios o financieros, ni
informaci\'on de terceros que permita identificarlos. Si durante la observaci\'on se llegara a ver
contenido personal de un equipo, no ser\'a capturado ni registrado.

\subsection*{3. Riesgos y c\'omo se reducen}
La participaci\'on no implica riesgos f\'isicos. Los riesgos posibles son sentirse inc\'omodo al
hablar de su trabajo o de problemas internos de la organizaci\'on, y que la informaci\'on se use de
forma indebida. Para reducirlos: puede omitir cualquier pregunta, pausar o terminar la actividad
cuando quiera; el equipo usar\'a el c\'odigo de participante indicado anteriormente en lugar de su
nombre; y ninguna opini\'on suya se comunicar\'a a terceros de forma que lo identifique.

\subsection*{4. Beneficios}
No recibir\'a pago ni beneficio econ\'omico. Su colaboraci\'on ayudar\'a a que la propuesta de
sistema responda a las necesidades reales de su trabajo y contribuir\'a a la formaci\'on de futuros
profesionales de la Ingenier\'ia de Software.

\subsection*{5. Confidencialidad y protecci\'on de sus datos}
Sus datos personales ser\'an tratados conforme a la Ley Org\'anica de Protecci\'on de Datos
Personales del Ecuador (LOPDP), que le reconoce, entre otros, los derechos de acceso,
rectificaci\'on, eliminaci\'on y oposici\'on sobre sus datos.
\begin{itemize}[nosep]
  \item \textbf{Lo que nunca se publica:} su nombre, n\'umero de c\'edula, firma, rostro, voz,
        contacto ni su cargo espec\'ifico. En todo documento, presentaci\'on o repositorio p\'ublico
        usted aparecer\'a solo con su c\'odigo de participante o con un rol gen\'erico (por ejemplo,
        ``Responsable del taller'').
  \item \textbf{D\'onde se guarda lo identificable:} este formulario firmado y las grabaciones
        originales se guardan en una carpeta institucional privada (cuenta \texttt{@uteq.edu.ec}),
        compartida solo con el docente responsable. En el repositorio p\'ublico del equipo
        (ISR401-ContX) no se sube nada que identifique a nadie: solo material anonimizado, el
        manifiesto de las evidencias privadas y una copia de este formulario con nombre, c\'edula,
        firma y n\'umero de tel\'efono cubiertos.
  \item \textbf{Cu\'anto tiempo:} el material identificable se conservar\'a como m\'aximo 24 meses
        desde el cierre del per\'iodo acad\'emico y luego ser\'a destruido. Solo se conservar\'an los
        instrumentos en blanco y los resultados totalmente anonimizados.
\end{itemize}

\subsection*{6. Participaci\'on voluntaria y derecho a retirarse}
Su participaci\'on es libre. Puede negarse a responder cualquier pregunta y retirarse en cualquier
momento, sin dar explicaciones y sin ninguna consecuencia laboral, acad\'emica ni de otro tipo. Si
se retira, puede pedir que se eliminen sus datos usando el desprendible de la \'ultima p\'agina o
escribiendo a los contactos indicados al inicio; se eliminar\'an siempre que a\'un no se hayan
integrado en resultados agregados o anonimizados de forma irreversible. Este formulario est\'a
dirigido a personas de 18 a\~nos o m\'as; las personas menores de edad no participan en el estudio.

\subsection*{7. Autorizaciones espec\'ificas}
Marque S\'i o No en cada fila. Puede aceptar unas y rechazar otras; solo la primera es necesaria
para participar.

\begin{tabularx}{\linewidth}{@{}Ycc@{}}
  \toprule
  \textbf{Autorizaci\'on} & \textbf{S\'i} & \textbf{No} \\
  \midrule
  a) Acepto participar en la actividad que se me ha explicado. & $\square$ & $\square$ \\
  b) Autorizo la grabaci\'on en audio de mi participaci\'on. & $\square$ & $\square$ \\
  c) Autorizo la grabaci\'on en video de mi participaci\'on. & $\square$ & $\square$ \\
  d) Autorizo fotograf\'ias fechadas del entorno de trabajo, sin mi rostro, sin documentos
     legibles, sin pantallas que muestren datos de clientes y sin ubicaci\'on GPS. & $\square$ & $\square$ \\
  e) Autorizo que se citen textualmente frases m\'ias, sin datos que me identifiquen. & $\square$ & $\square$ \\
  f) Autorizo el uso de la informaci\'on anonimizada en los trabajos acad\'emicos de la asignatura. & $\square$ & $\square$ \\
  \bottomrule
\end{tabularx}

\smallskip
Ninguna de estas autorizaciones permite publicar datos que lo identifiquen.

\subsection*{8. Declaraci\'on y firmas}
Declaro que tengo 18 a\~nos o m\'as; que he le\'ido este formulario, o me lo han le\'ido; que pude
hacer preguntas y recib\'i respuestas claras; y que comprendo el objetivo, las actividades, los
riesgos, los beneficios, el manejo de mis datos y mi derecho a retirarme. Acepto participar de forma
voluntaria seg\'un las opciones que marqu\'e en la Secci\'on 7. Se firman dos ejemplares: uno queda en
mi poder y otro en el expediente restringido del proyecto.

\vspace{1.2cm}
\noindent
\begin{tabular}{@{}p{7.4cm}@{\hspace{0.8cm}}p{7.4cm}@{}}
  \rule{7.2cm}{0.4pt} & \rule{7.2cm}{0.4pt} \\
  \centering Nombres y apellidos del participante & \centering\arraybackslash Firma del participante \\[1cm]
  \rule{7.2cm}{0.4pt} & \rule{7.2cm}{0.4pt} \\
  \centering N\'umero de c\'edula & \centering\arraybackslash Lugar y fecha (dd/mm/aaaa) \\[1cm]
  \rule{7.2cm}{0.4pt} & \rule{7.2cm}{0.4pt} \\
  \centering Nombre y firma del estudiante que aplica el instrumento &
  \centering\arraybackslash Nombre y firma del testigo (opcional; obligatorio si el formulario se ley\'o en voz alta) \\
\end{tabular}

\subsection*{9. Revocaci\'on del consentimiento (opcional) --- desprendible para el participante}
Conserve esta hoja. Si m\'as adelante desea retirarse del estudio, compl\'etela y entr\'eguela a
cualquier integrante del equipo o env\'ie su contenido por correo a los contactos indicados en la
primera p\'agina.

\medskip
\noindent Yo, \rule{6cm}{0.4pt}, participante con c\'odigo \rule{2.5cm}{0.4pt}, retiro mi
consentimiento para participar en el proyecto Sistema de gesti\'on de reparaci\'on de equipos de
c\'omputo y solicito:
\begin{itemize}[nosep]
  \item[$\square$] Eliminar mis datos identificables (formulario firmado, grabaciones, fotograf\'ias).
  \item[$\square$] Eliminar tambi\'en mis respuestas anonimizadas, siempre que a\'un no se hayan
                   integrado en resultados agregados.
\end{itemize}

\vspace{1.2cm}
\noindent
\begin{tabular}{@{}p{7.4cm}@{\hspace{0.8cm}}p{7.4cm}@{}}
  \rule{7.2cm}{0.4pt} & \rule{7.2cm}{0.4pt} \\
  \centering Firma & \centering\arraybackslash Fecha (dd/mm/aaaa) \\
\end{tabular}
